% !TeX root = MappingTheory.tex

\chapter{Логика}
\section{Утверждения}
Основой всей математики является логика, а именно доказательство различных утверждений, например 'Если число делится на 4, то оно делится на 2'. Любая отрасль математики по сути своей изначально строится на том, что мы берем несколько фактов, называемых аксиомами, которые принимаются за истину, формулируем множество определений, а затем с помощью логики выводим все интересующие нас утверждения. Под утверждением понимается любое заявление с которым мы хотим работать, тут нам необходимо только интуитивное понимание.
Цель этой главы -- познакомиться с основными принципами доказательств на примере логических утверждений и построить пример системы, где из нескольких аксиом и определений получаются все интересующие нас факты.

\quad \\
\underline{\textbf{Терминология}} \\

Любое утверждение с которым мы работаем -- мы рассматриваем, и мы \textbf{НЕ} говорим что оно верно или нет, т.к. это приводит к путанице. Вместо этого мы будем говорить что та или иная ситуация верна или ложна в нашем утверждении. Если под утверждением понимается универсальное заявление, то ситуация -- это любой вариант реальности, который мы можем себе представить.
\begin{example}
Утверждение 'У меня есть красные волосы', ситуация 'все волосы черные' ложна в этом утверждении, а ситуация 'у меня половина волос красная, а половина зеленая' уже верна.
\end{example}

Классически утверждения не прописывают каждый раз целиком, вместо этого их обозначают заглавными буквами латинского алфавита $ A,B,C \ldots $, но мы будем обозначать их курсивом $\mathcal{A}, \mathcal{M}, \mathcal{F} \ldots $, поскольку прямые буквы уже заняты другим объектом.
Ситуации же будем обозначать маленькими буквами латинского или греческого алфавита $ a, \omega, \xi \ldots $ \\

\begin{axiom}
Для любой ситуации и для всякого утверждения справедливо, что эта ситуация либо верна, либо ложна в этом утверждении.
В нашей системе не будет ситуаций одновременно верных и ложных в утверждении, мы запрещаем это первой аксиомой и если мы получили такой результат, то это противоречие.
\end{axiom}

Введем удобное сокращение, текст 'Ситуация $ \omega $ верна в утверждении $\mathcal{A}$' будем записывать как $\mathcal{A}(\omega)$.
Если вместе с этим словосочетание 'для всякого' и 'для любого' заменить знаком $ \forall $, то первая аксиома записывается очень компактно: \[
  \forall \mathcal{A}\; \forall \omega \quad \mathcal{A}(\omega) \text{ или } \omega \text{ ложна в } \mathcal{A} \]

\begin{definition}\label{negation}
  Пусть $\mathcal{A}$ -- утверждение, тогда отрицанием $\mathcal{A}$ будет называться утверждение, подходящее под следующие условия\\
  
  $ \forall $ (для всякой) ситуации $\omega$ справедливо: \begin{enumerate}
    \item Если $ \omega $ верна в $ \mathcal{A} $, то $\omega$ ложна в отрицании $\mathcal{A}$.
    \item Если $ \omega $ ложна в $\mathcal{A}$, то $ \omega $ верна в отрицании $\mathcal{A}$.
  \end{enumerate}
\end{definition}

Поскольку любая ситуация верна или ложна в $\mathcal{A}$ мы всегда можем узнать верна или ложна эта ситуация в отрицании $\mathcal{A}$. Отрицание $\mathcal{A}$  обозначается как $ \neg\mathcal{A} $ или $ \overline{\mathcal{A}} $.

\begin{exercise}
  Пусть утверждение $\mathcal{B}$ -- У меня нет волос. Сформулируйте $\neg\mathcal{B}$ и докажите что это 
  действительно отрицание, используя определение.
\end{exercise}

\begin{definition}
  Пусть $\mathcal{A}$ -- утверждение, $\mathcal{A}$ называется тавтологией, если \[
     \forall \omega \quad \mathcal{A}(\omega). \]
  Читается как, для всякой ситуации омега, омега верна в утверждении A.
\end{definition}
\begin{definition}
  Пусть $\mathcal{A}$ -- утверждение, $\mathcal{A}$ называется противоречием, если \[
     \forall \omega \quad \omega \text{ ложна в } \mathcal{A}. \]
\end{definition}

\begin{statement}
  Пусть $\mathcal{A}$ -- тавтология, тогда $ \neg \mathcal{A} $ -- противоречие.
\end{statement}
\begin{proof}
  Возьмем произвольную ситуацию $ \xi $, наша задача показать что $ \xi $ ложна в $ \neg \mathcal{A} $.
  По определению тавтологии мы знаем, что $ \mathcal{A}(\xi) $, а по определению отрицания \[
    \forall \omega \quad \text{ Если } \mathcal{A}(\omega) \text{, то } \omega \text{ ложна в } \neg\mathcal{A}. \]
  Отсюда делаем вывод, что $ \xi $ действительно ложна в $ \neg\mathcal{A} $. Здесь $ \xi $ -- абсолютно произвольная
  ситуация, а значит рассуждения справедливы для всех возможных ситуаций.
  Таким образом $ \neg\mathcal{A} $ -- противоречие.
\end{proof}
\begin{exercise}
  Докажите, что отрицание противоречия -- тавтология.
\end{exercise}
